\section{Plan de Modelamiento Estadístico}

Como ya se mencionó en la sección anterior, la observación principal en los gráficos es que la dependencia del colesterol a la edad es de forma cuadrática. Como esta relación no necesariamente está centrada en el cero, tampoco lo estará la hipotética parábola, consecuentemente, la relación estará dada necesariamente por dos covariables: edad y edad$^2$. Como  el objetivo es comparar con una relación estrictamente lineal, la pregunta de investigación se reduce a medir la importancia de la covariable edad$^2$. Considerando esto, se proponen dos modelos anidades con intercepto con las siguientes covariables: $X_1 = \texttt{RIAGENDIR} - 1$, sexo codificado en 0 para hombres y 1 para mujeres; $X_2 = \texttt{BMXBMI}$, índice de masa corportal; $X_3 = \texttt{INDFMPIR}$, razón de ingresos familiares respecto al umbral de pobreza, para medir el nivel socioeconómico del sujeto de la observación; $X_4 = \tfrac{\texttt{BPXOSY1} + \texttt{BPXOSY2} + \texttt{BPXOSY3}}{3}$, promedio simple de las tres mediciones de la presión sistólica; $X_5 = \tfrac{\texttt{BPXODI1} + \texttt{BPXODI2} + \texttt{BPXODI3}}{3}$, promedio simple de las tres mediciones de la presión diastólica; $X_6 = \texttt{RIDAGEYR}$, edad; $X_7 = \texttt{RIDAGEYR}^2$, edad al cuadrado. Estas covariables nos definen los siguientes modelos:

\begin{itemize}
    \item Modelo completo $\Omega$
    \begin{equation*}
        Y = \beta_0 + \beta_1 X_1 + \beta_2 X_2 + \beta_3 X_3 + \beta_4 X_4 + \beta_5 X_5 + \beta_6 X_6 + \beta_7 X_7 + \varepsilon
    \end{equation*}

    \item Modelo reducido $\omega$
    \begin{equation*}
        Y = \beta_0 + \beta_1 X_1 + \beta_2 X_2 + \beta_3 X_3 + \beta_4 X_4 + \beta_5 X_5 + \beta_6 X_6 + \varepsilon
    \end{equation*}
\end{itemize}

Donde $Y$ es la variable dependiente, en este caso, el colesterol. Con respecto al error $\varepsilon$ asumiremos $\mathbb{E}[\varepsilon | X] = 0$ exogeneidad media para que los estimadores por mínimos cuadrados sean insesgados, por otro lado, debido a que las observaciones son de personas independientes, es seguro asumir $\text{Cov}(\varepsilon | X) = \sigma^2 I_n$ homocedasticidad y no correlación condicional.

% TODO: Completar plan de modelamiento
% - Especificar modelos a utilizar (regresión lineal, splines, polinomios, etc.)
% - Detallar criterios de selección de modelo
% - Explicar validación cruzada y evaluación

% COMIENZA SECCIÓN DE DIAGNÓSTICOS

Teniendo claramente estblecido el modelo, a continuación se comentarán los diagnósticos previstos: en primera instancia se realizarán gráficos para conocer los datos y cómo resultó el ajuste. De los gráficos mencionados los primeros anáisis serán las proyecciones de los errores versus cada variable independiente, buscando detectear relaciones no lineales, heterocedasticidad o alguna combinación de estas. Además de los gráficos anteriores se realizará un un qqplot

Anteriormente se mencionó que no se integraron las variables que visualmente tenían relaciones directamente lineales para evitar problemas de colinealidad que empeoren el condicionamiento de la matriz de covariables. Para formalizar este análisis con las variables que si se integraron al modelo se revisará el condicionamiento de la matriz y los VIF por variable.

Los test más importante en cuanto a evaluar directamente nuestra pregunta de investigación es comparar modelos anidados incluyendo o no el coeficiente que acompaña al término de edad$^2$ (comparación de los modelos $\Omega$ y $\omega$ descritos anteriormente); además para ambos por separado se va a testear la posibilidad de reducir otras variables bajo criterio de eliminación, revisando valores de CP de Mallows y $R^2_a$.




